\begin{statement}[C.7.2; Proved: the compact-subgroup criterion and the real formula]
For $P$ as in Statement C.7.1, $N(P)$ is algebraic if and only
if the fibre of $G_P$ above $P$ contains an algebraic point of
the maximal compact subgroup $G_P(\C)^c$.
For the family in the prompt, if $t>1$ is real, choose
$\tau=iT$, $T>1$, and the zero $\mu=1/2-iy$, $0<y<T/2$.
With $x=2y/T$, one has
\[
 N(P)^4=3^{3/4}(t-1)^{-1/4}
       \left|\frac{\vartheta_1(x,i/T)}{\eta(i/T)}\right|.
\]
\end{statement}
\begin{proof}
On the analytic universal cover of $G_P$ define the real
homomorphism
\[
 \chi(z,v)=\Re\{v-\zeta(u)z+u\eta(z)\}.
\]
It vanishes on the periods \eqref{eq:cp-periods}, hence descends
to $G_P(\C)$. Its restriction to the torus is $\log|\cdot|$.
The kernel is compact: in the universal cover it is a real
three-dimensional vector space modulo the rank-three period
lattice. Every compact subgroup has zero image in $\R$, so
this kernel is exactly $G_P(\C)^c$.
The identity
$\Re(u\eta(z))=\Re(z\eta(u))$ follows from the Legendre relation;
equivalently $\eta(z)=sz+\pi\bar z/\operatorname{area}(\Lambda)$
for a complex constant $s$. Substituting
\eqref{eq:cp-v} at $z=u$ gives
\begin{equation}\label{eq:cp-chi}
 \chi(R)=\log|r|-\log|\wp'(u)|
            -2\log N(P)+\frac16\log|\Delta|.
\end{equation}
If $r$ is algebraic and $\chi(R)=0$, this equation makes
$N(P)^2$ algebraic. Conversely, if $N(P)$ is algebraic,
choose the positive real algebraic number
$r=|\wp'(u)|\,|\Delta|^{-1/6}N(P)^2$.
Then \eqref{eq:cp-chi} vanishes. This proves both directions
with the chosen rational trivialization of the fibre.
The compact interpretation is the one discussed by
\citet[slides 5--8]{bertrand2017}; here the calculation fixes
its normalization on the diagonal $P=P$.

For the real formula, the rectangular elliptic curve has its
two zeros of $\wp$ on $1/2+i\R$; along that edge $\wp$ runs
between the positive root $e_1$ and the negative root $e_2$.
Thus the indicated representative exists. Translate $2\mu$
by $-1$ and apply the modular substitution $\tau\mapsto-1/\tau$.
The resulting elliptic argument is $-2y/T$, which is real.
The Gaussian is therefore one. Evenness of $N$ and the
duplication identity prove the displayed formula. This removes
the Gaussian from this particular expression but makes no
algebraicity assertion about its theta value.
\end{proof}
